% GENERATED FILE. Edit canonical lessons, then run npm run generate:books.
% canonical-source-hash: fnv1a64:ead26caf2abd2b3b
% canonical-lessons: TE-C155-vacce, TE-C155-ikapai, TE-C155-bhavishyattu, TE-C155-ashincu, TE-C155-edurucudu

\chapter{Next, From Now On, the Future, To Hope, To Look Forward To}
\label{ch:te-next-from-now-on-the-future-to-hope-to-look-forward-to}
\hlchaptermodality{155}

\begin{chapteropening}
\textbf{By the end of this chapter:} \emph{I can say what will happen, and when: next week, from now on, and what I hope for.}

Complete the last lesson of chapter 155: I can say what will happen, and when: next week, from now on, and what I hope for.
\end{chapteropening}

\section[\texorpdfstring{\te{వచ్చే} (vaccē)}{వచ్చే (vaccē)}]{\te{వచ్చే} (vaccē) — next, coming}

\label{lesson:TE-C155-vacce}



\begin{quote}
Before the new one: say the Telugu for once, then the Telugu for to spend.
\end{quote}

\subsection*{You'll want to know: \te{వచ్చే}}
\textbf{\te{వచ్చే}} — \emph{vaccē} — \textquotedblleft{}next, coming\textquotedblright{}. The future for \textquotedblleft{}I\textquotedblright{} ends in \textbf{-\te{తాను}} (\emph{-tānu}): \textbf{\te{వచ్చే} \te{వారం} \te{నేను} \te{విజయవాడ} \te{వెళ్తాను}} — \emph{vaccē vāraṁ nēnu vijayavāḍa veḷtānu} — \textquotedblleft{}Next week I will go to Vijayawada.\textquotedblright{}

\begin{grammarlens}[title={What you've built}]
One more word for this chapter.
\end{grammarlens}

\subsection*{Your turn}
\begin{itemize}
\raggedright
  \item \emph{Say it:} \emph{vaccē}
  \item \emph{Say it:} \emph{vaccē}, once more
  \item \emph{Say it:} say \emph{vaccē}
  \item \emph{Recall:} say the Telugu for once, then the Telugu for to spend, then say \emph{vaccē} again
\end{itemize}

\begin{tcolorbox}[breakable,colback=teal!4,colframe=teal!35!black,title={Before you move on}]
What does \te{వచ్చే} mean? (\textquotedblleft{}Next, coming\textquotedblright{}.) Say it once more.
\end{tcolorbox}
\section[\texorpdfstring{\te{ఇకపై} (ikapai)}{ఇకపై (ikapai)}]{\te{ఇకపై} (ikapai) — from now on}

\label{lesson:TE-C155-ikapai}



\begin{quote}
Before the new one: say the Telugu for to spend, then the Telugu for next.
\end{quote}

\subsection*{You'll want to know: \te{ఇకపై}}
\textbf{\te{ఇకపై}} — \emph{ikapai} — \textquotedblleft{}from now on\textquotedblright{}. \textbf{\te{ఇకపై} \te{నేను} \te{త్వరగా} \te{లేస్తాను}} — \emph{ikapai nēnu tvaragā lēstānu} — \textquotedblleft{}From now on I will get up early.\textquotedblright{}

\begin{grammarlens}[title={What you've built}]
One more word for this chapter.
\end{grammarlens}

\subsection*{Your turn}
\begin{itemize}
\raggedright
  \item \emph{Say it:} \emph{ikapai}
  \item \emph{Say it:} \emph{ikapai}, once more
  \item \emph{Say it:} say \emph{ikapai}
  \item \emph{Recall:} say the Telugu for to spend, then the Telugu for next, then say \emph{ikapai} again
\end{itemize}

\begin{tcolorbox}[breakable,colback=teal!4,colframe=teal!35!black,title={Before you move on}]
What does \te{ఇకపై} mean? (\textquotedblleft{}From now on\textquotedblright{}.) Say it once more.
\end{tcolorbox}
\section[\texorpdfstring{\te{భవిష్యత్తు} (bhaviṣyattu)}{భవిష్యత్తు (bhaviṣyattu)}]{\te{భవిష్యత్తు} (bhaviṣyattu) — the future}

\label{lesson:TE-C155-bhavishyattu}



\begin{quote}
Before the new one: say the Telugu for next, then the Telugu for from now on.
\end{quote}

\subsection*{You'll want to know: \te{భవిష్యత్తు}}
\textbf{\te{భవిష్యత్తు}} — \emph{bhaviṣyattu} — \textquotedblleft{}the future\textquotedblright{}. \textbf{\te{భవిష్యత్తులో} \te{నేను} \te{ఉపాధ్యాయుడిని} \te{అవుతాను}} — \emph{bhaviṣyattulō nēnu upādhyāyuḍini avutānu} — \textquotedblleft{}In the future I will become a teacher.\textquotedblright{}

\begin{grammarlens}[title={What you've built}]
One more word for this chapter.
\end{grammarlens}

\subsection*{Your turn}
\begin{itemize}
\raggedright
  \item \emph{Say it:} \emph{bhaviṣyattu}
  \item \emph{Say it:} \emph{bhaviṣyattu}, once more
  \item \emph{Say it:} say \emph{bhaviṣyattu}
  \item \emph{Recall:} say the Telugu for next, then the Telugu for from now on, then say \emph{bhaviṣyattu} again
\end{itemize}

\begin{tcolorbox}[breakable,colback=teal!4,colframe=teal!35!black,title={Before you move on}]
What does \te{భవిష్యత్తు} mean? (\textquotedblleft{}The future\textquotedblright{}.) Say it once more.
\end{tcolorbox}
\section[\texorpdfstring{\te{ఆశించు} (āśiñcu)}{ఆశించు (āśiñcu)}]{\te{ఆశించు} (āśiñcu) — to hope, to expect}

\label{lesson:TE-C155-ashincu}



\begin{quote}
Before the new one: say the Telugu for from now on, then the Telugu for the future.
\end{quote}

\subsection*{You'll want to know: \te{ఆశించు}}
\textbf{\te{ఆశించు}} — \emph{āśiñcu} — \textquotedblleft{}to hope, to expect\textquotedblright{}. \textbf{\te{మీరు} \te{బాగున్నారని} \te{ఆశిస్తున్నాను}} — \emph{mīru bāgunnārani āśistunnānu} — \textquotedblleft{}I hope you are well.\textquotedblright{}

\begin{grammarlens}[title={What you've built}]
One more word for this chapter.
\end{grammarlens}

\subsection*{Your turn}
\begin{itemize}
\raggedright
  \item \emph{Say it:} \emph{āśiñcu}
  \item \emph{Say it:} \emph{āśiñcu}, once more
  \item \emph{Say it:} say \emph{āśiñcu}
  \item \emph{Recall:} say the Telugu for from now on, then the Telugu for the future, then say \emph{āśiñcu} again
\end{itemize}

\begin{tcolorbox}[breakable,colback=teal!4,colframe=teal!35!black,title={Before you move on}]
What does \te{ఆశించు} mean? (\textquotedblleft{}To hope, to expect\textquotedblright{}.) Say it once more.
\end{tcolorbox}
\section[\texorpdfstring{\te{ఎదురుచూడు} (edurucūḍu)}{ఎదురుచూడు (edurucūḍu)}]{\te{ఎదురుచూడు} (edurucūḍu) — to look forward to, to wait for}

\label{lesson:TE-C155-edurucudu}



\begin{quote}
Before the new one: say the Telugu for the future, then the Telugu for to hope.
\end{quote}

\subsection*{You'll want to know: \te{ఎదురుచూడు}}
\textbf{\te{ఎదురుచూడు}} — \emph{edurucūḍu} — \textquotedblleft{}to look forward to, to wait for\textquotedblright{}. \textbf{\te{నేను} \te{మీ} \te{ఉత్తరం} \te{కోసం} \te{ఎదురుచూస్తున్నాను}} — \emph{nēnu mī uttaraṁ kōsaṁ edurucūstunnānu} — \textquotedblleft{}I am looking forward to your letter.\textquotedblright{}

\begin{grammarlens}[title={What you've built}]
One more word for this chapter.
\end{grammarlens}

\subsection*{Your turn}
\begin{itemize}
\raggedright
  \item \emph{Say it:} \emph{edurucūḍu}
  \item \emph{Say it:} \emph{edurucūḍu}, once more
  \item \emph{Say it:} say \emph{edurucūḍu}
  \item \emph{Recall:} say the Telugu for the future, then the Telugu for to hope, then say \emph{edurucūḍu} again
\end{itemize}

\begin{tcolorbox}[breakable,colback=teal!4,colframe=teal!35!black,title={Before you move on}]
What does \te{ఎదురుచూడు} mean? (\textquotedblleft{}To look forward to, to wait for\textquotedblright{}.) Say it once more.
\end{tcolorbox}
